% Generated draft: requires mathematical and source review.
\documentclass[11pt]{article}
\usepackage[margin=25mm]{geometry}
\usepackage{fontspec}
\setmainfont{Latin Modern Roman}
\usepackage{amsmath,amssymb,mathtools,unicode-math}
\setmathfont{Latin Modern Math}
\usepackage{enumitem,needspace,xurl,hyperref}
\setlength{\parindent}{0pt}
\setlength{\parskip}{5pt}
\setlength{\emergencystretch}{3em}
\newcommand{\N}{\mathbb{N}}
\newcommand{\Z}{\mathbb{Z}}
\newcommand{\Q}{\mathbb{Q}}
\newcommand{\R}{\mathbb{R}}
\newcommand{\C}{\mathbb{C}}
\newcommand{\F}{\mathbb{F}}
\newcommand{\A}{\mathbb{A}}
\newcommand{\PP}{\mathbb{P}}
\newcommand{\E}{\mathbb{E}}
\newcommand{\eps}{\varepsilon}
\newcommand{\dd}{\,\mathrm{d}}
\begin{document}
\section*{UnsolvedMath ID: 30001413}
\textbf{Source problem code:} catalogue entry
\subsection{Definitions and mathematical statement}
Let \(\Omega\) denote the domain of the hinged biharmonic plate problem, and let \(v\) and \(f\) denote the displacement and forcing, respectively. Write \(\Delta\) for the Laplacian and \(\Delta^2=\Delta\circ\Delta\). The supplied boundary value problem is
\[
\begin{cases}
\Delta^2v=f & \text{in }\Omega,\\
v=0 & \text{on }\partial\Omega,\\
\Delta v-(1-\sigma)\kappa v_n=0 & \text{on }\partial\Omega.
\end{cases}
\]
Here \(\sigma\) is the parameter appearing in the boundary condition, \(\kappa\) denotes boundary curvature, and \(v_n\) denotes a normal derivative. The supplied record does not specify the range of \(\sigma\), the ambient dimension, the curvature convention, or the orientation of the normal derivative; these conventions must be fixed together before the boundary condition is fully determined.

A boundary is of class \(C^{2,1}\) if it is locally representable by graphs of twice continuously differentiable functions whose second derivatives are locally Lipschitz continuous. The positivity property under discussion is
\[
 f\geq 0\text{ in }\Omega\quad\Longrightarrow\quad v\geq 0\text{ in }\Omega,
\]
for every admissible forcing \(f\) and its corresponding solution \(v\) of the displayed problem. The admissible function spaces, solution concept, and pointwise or almost-everywhere interpretation of these inequalities are not supplied; neither are existence and uniqueness hypotheses.

\textbf{Problem.} The source record describes convexity of \(\Omega\) together with \(C^{2,1}\) regularity of \(\partial\Omega\) as sufficient domain assumptions for this positivity property. Determine whether these domain assumptions are optimal or whether they can be weakened while preserving the positivity property for every admissible nonnegative forcing. This is the supplied research question, conditional on the omitted parameter and analytic hypotheses of the underlying result. The record does not specify a class of weaker domains or a precise meaning of optimality, so it does not determine a single fully quantified sharpness assertion.
\subsection{Short English statement}
For the stated hinged plate boundary conditions, determine whether positivity of the displacement for every nonnegative load persists under domain assumptions weaker than convexity and a \(C^{2,1}\) boundary. The parameter range, curvature conventions, analytic setting, and intended meaning of optimality require clarification from the underlying source.
\subsection{Sources}
\begin{itemize}
\item Supplied local UnsolvedMath registry record, permanent ID \(30001413\), title: \emph{Optimal Domain Conditions for Positivity of Hinged Plates}; exact source problem code: \texttt{catalogue entry}. The record assigns the problem to 2010 and supplies no individual proposer.
\item Supplied report citation: \emph{Optimal Constants in the Theory of Sobolev Spaces and PDEs} (2010), Oberwolfach Reports, pp.\ 325--371, DOI 10.4171/owr/2010/08. \url{https://doi.org/10.4171/owr/2010/08}. The background identifies Oberwolfach Reports as reports of mathematical workshops and their open-problem discussions.
\item Supplied literature source, labelled ``Oberwolfach Report 8/2010'': \url{https://doi.org/10.4171/owr/2010/08}.
\item Supplied literature assessment checked on \texttt{2026-08-21}: the record labels the problem open, reports that the report leaves optimality of the domain assumptions open, and states that the later literature located in that curation does not establish the requested positivity principle under weaker general assumptions. This is supplied dated curation, not an independently verified literature assessment. The supplied background expressly requests independent verification of time-sensitive status and cited claims.
\end{itemize}
\end{document}
